\documentclass[11pt,a4paper]{article}
\usepackage[utf8]{inputenc}
\usepackage[T1]{fontenc}
\usepackage{amsmath, amssymb, amsthm}
\usepackage{geometry}
\geometry{margin=1in}
\usepackage{fancyhdr}
\usepackage{listings}
\usepackage{newunicodechar}
\usepackage{hyperref}

\pagestyle{fancy}
\fancyhf{}
\rhead{Erd\H{o}s-Straus Conjecture}
\cfoot{\thepage}

\title{Rigorous Formulation and Proof Strategy for the Erd\H{o}s-Straus Conjecture}
\author{Charles EDOU NZE, chercheur ind\'ependant}
\date{\today}

\newtheorem{theorem}{Theorem}[section]
\newtheorem{lemma}[theorem]{Lemma}
\newtheorem{definition}[theorem]{Definition}
\newtheorem{proposition}[theorem]{Proposition}

\newunicodechar{ℕ}{\ensuremath{\mathbb{N}}}
\newunicodechar{∀}{\ensuremath{\forall}}
\newunicodechar{∃}{\ensuremath{\exists}}
\newunicodechar{∧}{\ensuremath{\wedge}}
\newunicodechar{∨}{\ensuremath{\vee}}
\newunicodechar{→}{\ensuremath{\rightarrow}}

\begin{document}

\maketitle

\begin{abstract}
This document presents a rigorous axiomatic foundation, literature context, modular identity strategies, and an explicit mathematical derivation concerning the Erd\H{o}s-Straus conjecture. Additionally, a structural architecture for autoformalization in Lean 4 is provided.

\vspace{2cm}
\begin{flushright}
\textit{Charles EDOU NZE, chercheur ind\'ependant}
\end{flushright}
\end{abstract}

\newpage

\section{Axiomatic Definitions}

\begin{definition}
Let $\mathbb{N}_{>1} = \{n \in \mathbb{N} \mid n \geq 2\}$. The Erd\H{o}s-Straus conjecture postulates that for all $n \in \mathbb{N}_{>1}$, there exist $x, y, z \in \mathbb{N}_{>0}$ such that:
\begin{equation}
\frac{4}{n} = \frac{1}{x} + \frac{1}{y} + \frac{1}{z}
\end{equation}
\end{definition}

\begin{definition}
By multiplying the aforementioned equation by $n x y z$, we obtain the equivalent polynomial Diophantine formulation:
\begin{equation}
4 x y z = n (x y + x z + y z)
\end{equation}
\end{definition}

\section{Contextual Literature and Related Works}

The Erd\H{o}s-Straus conjecture, formulated in 1948 by Paul Erd\H{o}s and Ernst G. Straus, seeks an expansion of $4/n$ into exactly three positive Egyptian fractions.
Relevant literature extensively explores modular identities. For instance, Mordell (1967) demonstrated that an identity providing solutions for all $n \equiv r \pmod p$ exists if and only if $r$ is not a quadratic residue modulo $p$.
Computer searches (e.g., by Salez, 2014) have verified the conjecture up to $N = 10^{17}$.
Elsholtz and Tao (2013) bounded the average number of solutions polylogarithmically using the Bombieri-Vinogradov and Brun-Titchmarsh theorems.

The Hasse principle suggests that studying Diophantine equations modulo $p$ might piece together integer solutions. However, a complete covering system of congruences is proven impossible due to 1 being a quadratic residue modulo any $n > 1$.

\section{Strategy of Proof and Lemmas}

The problem can be reduced to considering only prime values of $n$. If $n = m \cdot p$ where $p$ is prime, a solution for $p$ yields a solution for $n$ by scaling the denominators by $m$.

\begin{lemma}[Modular Case $n \equiv 2 \pmod 3$]
For any integer $n > 1$ such that $n \equiv 2 \pmod 3$, there exists a solution $(x,y,z) \in \mathbb{N}_{>0}^3$ satisfying $\frac{4}{n} = \frac{1}{x} + \frac{1}{y} + \frac{1}{z}$.
\end{lemma}

\begin{lemma}[Modular Case $n \equiv 3 \pmod 4$]
For any integer $n > 1$ such that $n \equiv 3 \pmod 4$, there exists a solution $(x,y,z) \in \mathbb{N}_{>0}^3$ satisfying $\frac{4}{n} = \frac{1}{x} + \frac{1}{y} + \frac{1}{z}$.
\end{lemma}

\section{Informal Mathematical Proof (Zero Ellipse)}

We proceed to prove the first lemma explicitly.

\begin{proof}[Proof of Lemma 3.1]
Let $n \in \mathbb{N}_{>1}$ be such that $n \equiv 2 \pmod 3$.
By definition of congruence modulo 3, there exists an integer $k \geq 0$ such that $n = 3k + 2$.
We define $x = n$. Since $n \geq 2$, $x \in \mathbb{N}_{>0}$.
We define $y = \frac{n+1}{3}$.
We must verify that $y$ is an integer.
Substitute $n = 3k + 2$ into the expression for $y$:
\begin{align*}
y &= \frac{(3k + 2) + 1}{3} \\
  &= \frac{3k + 3}{3} \\
  &= k + 1
\end{align*}
Since $k \geq 0$, $k+1$ is an integer and $y \in \mathbb{N}_{>0}$.
We define $z = \frac{n(n+1)}{3}$.
We must verify that $z$ is an integer.
Substitute $n = 3k + 2$:
\begin{align*}
z &= \frac{(3k+2)(3k+3)}{3} \\
  &= (3k+2)(k+1)
\end{align*}
Since $k \geq 0$, $3k+2 \geq 2$ and $k+1 \geq 1$. Both are integers, so their product is an integer, and $z \in \mathbb{N}_{>0}$.

We now compute the sum $\frac{1}{x} + \frac{1}{y} + \frac{1}{z}$:
\begin{align*}
\frac{1}{x} + \frac{1}{y} + \frac{1}{z} &= \frac{1}{n} + \frac{1}{\frac{n+1}{3}} + \frac{1}{\frac{n(n+1)}{3}} \\
&= \frac{1}{n} + \frac{3}{n+1} + \frac{3}{n(n+1)}
\end{align*}
To add these fractions, we find a common denominator, which is $n(n+1)$.
\begin{align*}
\frac{1}{n} &= \frac{n+1}{n(n+1)} \\
\frac{3}{n+1} &= \frac{3n}{n(n+1)} \\
\frac{3}{n(n+1)} &= \frac{3}{n(n+1)}
\end{align*}
Summing the numerators:
\begin{align*}
\frac{n+1 + 3n + 3}{n(n+1)} &= \frac{4n + 4}{n(n+1)} \\
&= \frac{4(n+1)}{n(n+1)}
\end{align*}
We divide the numerator and the denominator by $(n+1)$. Since $n \geq 2$, $n+1 \neq 0$.
\begin{align*}
\frac{4(n+1)}{n(n+1)} = \frac{4}{n}
\end{align*}
Thus, $\frac{1}{x} + \frac{1}{y} + \frac{1}{z} = \frac{4}{n}$.
This completes the proof.
\end{proof}

\begin{proof}[Proof of Lemma 3.2]
Let $n \in \mathbb{N}_{>1}$ be such that $n \equiv 3 \pmod 4$.
By definition, there exists an integer $k \geq 0$ such that $n = 4k + 3$.
We define $x = \frac{n+1}{4}$.
Substitute $n = 4k + 3$:
\begin{align*}
x &= \frac{(4k+3)+1}{4} \\
  &= \frac{4k+4}{4} \\
  &= k+1
\end{align*}
Since $k \geq 0$, $x \in \mathbb{N}_{>0}$.
We define $y = \frac{n(n+1)}{2}$.
Substitute $n = 4k + 3$:
\begin{align*}
y &= \frac{(4k+3)(4k+4)}{2} \\
  &= \frac{(4k+3)4(k+1)}{2} \\
  &= 2(4k+3)(k+1)
\end{align*}
Since $k \geq 0$, $4k+3 \geq 3$ and $k+1 \geq 1$. Both are integers, so their product multiplied by 2 is an integer, and $y \in \mathbb{N}_{>0}$.
We define $z = \frac{n(n+1)}{2}$.
Since this expression is identical to $y$, $z \in \mathbb{N}_{>0}$.

We now compute the sum $\frac{1}{x} + \frac{1}{y} + \frac{1}{z}$:
\begin{align*}
\frac{1}{x} + \frac{1}{y} + \frac{1}{z} &= \frac{1}{\frac{n+1}{4}} + \frac{1}{\frac{n(n+1)}{2}} + \frac{1}{\frac{n(n+1)}{2}} \\
&= \frac{4}{n+1} + \frac{2}{n(n+1)} + \frac{2}{n(n+1)}
\end{align*}
Summing the last two fractions:
\begin{align*}
\frac{2}{n(n+1)} + \frac{2}{n(n+1)} &= \frac{4}{n(n+1)}
\end{align*}
Now, sum with the first fraction:
\begin{align*}
\frac{4}{n+1} + \frac{4}{n(n+1)}
\end{align*}
Find the common denominator $n(n+1)$:
\begin{align*}
\frac{4}{n+1} &= \frac{4n}{n(n+1)}
\end{align*}
Summing the numerators:
\begin{align*}
\frac{4n + 4}{n(n+1)} &= \frac{4(n+1)}{n(n+1)}
\end{align*}
Divide the numerator and denominator by $(n+1)$. Since $n \geq 2$, $n+1 \neq 0$:
\begin{align*}
\frac{4(n+1)}{n(n+1)} = \frac{4}{n}
\end{align*}
Thus, $\frac{1}{x} + \frac{1}{y} + \frac{1}{z} = \frac{4}{n}$.
This completes the proof.
\end{proof}

\section{Architecture for Autoformalization in Lean 4}

To facilitate automated verification by a formal system such as Lean 4, the proofs can be codified as follows.

\begin{lstlisting}[basicstyle=\ttfamily\small, breaklines=true]
import Mathlib.Data.Nat.Basic
import Mathlib.Data.Nat.Parity

/--
The Erdos-Straus Conjecture
-/
def erdos_straus_conjecture : Prop :=
  ∀ n : ℕ, n > 1 → ∃ x y z : ℕ, x > 0 ∧ y > 0 ∧ z > 0 ∧
    4 * x * y * z = n * (x * y + x * z + y * z)

/--
Lemma 1: Solution for n = 2 mod 3
-/
theorem erdos_straus_mod3_2 (n : ℕ) (h1 : n > 1) (h2 : n % 3 = 2) :
  ∃ x y z : ℕ, x > 0 ∧ y > 0 ∧ z > 0 ∧
    4 * x * y * z = n * (x * y + x * z + y * z) := by
  admit

/--
Lemma 2: Solution for n = 3 mod 4
-/
theorem erdos_straus_mod4_3 (n : ℕ) (h1 : n > 1) (h2 : n % 4 = 3) :
  ∃ x y z : ℕ, x > 0 ∧ y > 0 ∧ z > 0 ∧
    4 * x * y * z = n * (x * y + x * z + y * z) := by
  admit
\end{lstlisting}

\end{document}
